\subsection*{综合例题前的方法整理}

下面按题型归纳前面的工具，着重区分不同模型的使用条件。

\paragraph{Lorentz 力与匀强磁场轨道}

带电粒子在电磁场中受
\[
 \vec F=q(\vec E+\vec v\times\vec B).
\]
磁力始终垂直于瞬时速度，功率 $\vec F_B\cdot\vec v=0$，所以静磁场只改变速度方向，不改变速率
和动能。方向必须同时考虑外积和电荷符号：先按正电荷用右手定则判断，再对负电荷反向。

把速度分解为平行与垂直磁场的两部分，$v_\parallel$ 保持不变，$v_\perp$ 作圆周运动：
\[
 R=\frac{mv_\perp}{|q|B},\qquad
 \omega_c=\frac{|q|B}{m},\qquad
 T=\frac{2\pi m}{|q|B}.
\]
合运动是螺旋线，螺距为 $h=v_\parallel T$。非相对论情形下周期与速率无关，这是回旋加速器能用
固定频率交变电场反复加速的基础；当速率接近光速，等效惯性增大，粒子与固定频率逐渐失谐。

有限磁场区域的题目还要由初始点和速度方向确定圆心，再求轨道圆与边界的交点。只使用半径公式
而不画圆心位置，容易选错圆弧和出射方向。自聚焦装置中，不同初始角的粒子能在特定几何条件下
重新会聚，原因是同种粒子具有相同回旋频率，近轴条件下到达屏的时间又近似相同；
它们的 $p_\perp=p\sin\theta$ 随出射角改变，回旋半径并不相同。

\paragraph{漂移、选择器与 Hall 场}

互相垂直且均匀的电、磁场中，粒子回旋中心以
\[
 \vec v_D=\frac{\vec E\times\vec B}{B^2}
\]
漂移。这个 $\vec E\times\vec B$ 漂移与粒子质量、电量大小及正负无关；回旋方向仍由电荷符号决定。
若装置要求粒子不偏转，则 $q\vec E+q\vec v\times\vec B=0$，可筛选满足
$v=E/B$ 的速度分量，这就是速度选择器。两种条件看似相近，一个描述导向中心运动，一个描述给定
直线轨道，使用时要区分。

载流导体置于磁场中，载流子横向偏转并在两侧积累电荷，直到横向电场力抵消磁力。Hall 电压的
极性能够判断主要载流子符号，大小则给出载流子浓度。若电流、磁场和样品厚度分别为 $I,B,d$，
取常规电流沿 $+x$、磁场沿 $+z$、横向两侧为 $y=\pm w/2$，并定义
$U_H=V(-w/2)-V(w/2)$，单一载流子模型给出
\[
 U_H=\frac{IB}{nqd}.
\]
Corbino 圆环没有普通样品那样的横向边界积累，径向电流在磁场中产生环向响应，适合从另一几何
角度理解磁阻与 Hall 动力学。

\paragraph{磁矩、Zeeman 效应与阻尼运动}

电流回路的磁矩为 $\vec\mu=I\vec S$，在匀强磁场中受力矩
$\vec M=\vec\mu\times\vec B$，势能 $U=-\vec\mu\cdot\vec B$。运动点电荷的闭合圆轨道也等效为电流，
其轨道磁矩与角动量满足
\[
 \vec\mu=\frac{q}{2m}\vec L.
\]
外磁场使不同磁矩取向具有不同能量；正常 Zeeman 效应的经典图像可由轨道进动理解，而精确谱线
结构还需要量子化和电子自旋。

若同时存在与速度成正比、方向相反的阻力 $-m\gamma\vec v$，磁场运动成为半径逐渐减小的螺旋。
用复变量 $u=v_x+iv_y$ 可把两个耦合方程合为
$\dot u=-(\gamma+i\Omega)u$，故速率按 $e^{-\gamma t}$ 衰减，方向仍以回旋频率转动。加入恒定电场后，
解可分成衰减的暂态和稳定漂移；先求长期极限再叠加暂态，往往比直接解坐标分量更清楚。
